\section*{Bab \ref{ch:g7:chemical-reactions} --- Reaksi Kimia: Pereaksi dan Produk}

\begin{solution}{exo:g7:chemical-reactions:1}
Fisika: es yang mencair, gula yang \omterm{def:g2:mixing-and-dissolving:dissolve}{larut}. Kimia: korek api yang
\omterm{def:g4:what-a-fire-needs:burning}{terbakar}, roti yang dipanggang, paku yang berkarat.
\end{solution}

\begin{solution}{exo:g7:chemical-reactions:2}
\omterm{def:g7:chemical-reactions:reactant}{Pereaksi} adalah spesies yang terpakai; \omterm{def:g7:chemical-reactions:reactant}{produk} adalah spesies baru yang
terbentuk.
\end{solution}

\begin{solution}{exo:g7:chemical-reactions:3}
karbon $+$ oksigen $\longrightarrow$ karbon dioksida.
\end{solution}

\begin{solution}{exo:g7:chemical-reactions:4}
\omterm{def:g7:chemical-reactions:reactant}{Pereaksi}: besi dan oksigen. \omterm{def:g7:chemical-reactions:reactant}{Produk}: oksida besi.
\end{solution}

\begin{solution}{exo:g7:chemical-reactions:5}
Embunnya menunjukkan air; air kapur yang keruh menunjukkan karbon
dioksida.
\end{solution}

\begin{solution}{exo:g7:chemical-reactions:6}
seng $+$ asam klorida $\longrightarrow$ seng klorida $+$ hidrogen.
\end{solution}

\begin{solution}{exo:g7:chemical-reactions:7}
Sebelum: 4 (dalam \omterm{def:g7:atoms-and-molecules:molecule}{molekul} metana). Sesudah: $2 + 2 = 4$ (dua dalam
setiap \omterm{def:g7:atoms-and-molecules:molecule}{molekul} air).
\end{solution}

\begin{solution}{exo:g7:chemical-reactions:8}
Tidak ada spesies baru yang muncul: garamnya masih ada, tersebar di
seluruh air, dan dapat diperoleh kembali dengan \omterm{def:g3:separating-mixtures:evaporate}{menguapkan} airnya. Itu
\omterm{def:g7:chemical-reactions:physical-transformation}{perubahan fisika}.
\end{solution}

\begin{solution}{exo:g7:chemical-reactions:9}
Sebagian besar kayu itu telah bereaksi dengan oksigen menghasilkan
\omterm{def:g7:chemical-reactions:reactant}{produk-produk} berupa gas, karbon dioksida dan uap air, yang pergi ke
\omterm{def:g4:air-a-mixture-of-gases:air}{udara}. Hanya abunya yang tertinggal.
\end{solution}

\begin{solution}{exo:g7:chemical-reactions:10}
hidrogen $+$ oksigen $\longrightarrow$ air. \omterm{def:g7:chemical-reactions:reactant}{Produknya} membirukan
tembaga sulfat anhidrat.
\end{solution}

\begin{solution}{exo:g7:chemical-reactions:11}
Sebelum: $2 \times 2 = 4$ \omterm{def:g7:atoms-and-molecules:atom}{atom} hidrogen dan 2 \omterm{def:g7:atoms-and-molecules:atom}{atom} oksigen. Sesudah: dua
\omterm{def:g7:atoms-and-molecules:molecule}{molekul} air mengandung $2 \times 2 = 4$ \omterm{def:g7:atoms-and-molecules:atom}{atom} hidrogen dan $2 \times 1 =
2$ \omterm{def:g7:atoms-and-molecules:atom}{atom} oksigen. \omterm{def:g7:atoms-and-molecules:atom}{Atom-atom} yang sama, dalam jumlah yang sama.
\end{solution}

\begin{solution}{exo:g7:chemical-reactions:12}
\omterm{def:g7:atoms-and-molecules:atom}{Atom} oksigen tidak dapat muncul entah dari mana: gambar itu memiliki 2
\omterm{def:g7:atoms-and-molecules:atom}{atom} oksigen sebelum dan 3 sesudahnya. Yang benar: satu \omterm{def:g7:atoms-and-molecules:atom}{atom} karbon $+$
satu \omterm{def:g7:atoms-and-molecules:molecule}{molekul} oksigen $\longrightarrow$ satu \omterm{def:g7:atoms-and-molecules:molecule}{molekul} karbon dioksida,
tanpa sisa apa pun.
\end{solution}

\begin{solution}{pb:g7:chemical-reactions:1}
\textbf{1.} Metana dan oksigen.

\textbf{2.} Air.

\textbf{3.} Karbon dioksida.

\textbf{4.} Tidak: nitrogen tidak terpakai, dan tidak dituliskan dalam
\omterm{def:g7:chemical-reactions:word-equation}{persamaan kata}.

\textbf{5.} metana $+$ oksigen $\longrightarrow$ karbon dioksida $+$
air.

\textbf{6.} Kimia: dua spesies hilang (metana, oksigen), dan dua spesies
baru muncul, yang dideteksi dengan uji; selain itu panas dan cahaya
dilepaskan.

\textbf{7.} \ce{CH4}, \ce{O2}, \ce{CO2}, \ce{H2O}.

\textbf{8.} 1 karbon, 4 hidrogen, $2 \times 2 = 4$ oksigen.

\textbf{9.} 1 karbon; $2 \times 2 = 4$ hidrogen; $2 + 2 \times 1 = 4$
oksigen. \omterm{def:g7:atoms-and-molecules:atom}{Atom-atom} yang sama dalam jumlah yang sama: \omterm{def:g7:atoms-and-molecules:atom}{atom-atom} itu hanya
ditata ulang.

\textbf{10.} $10 \times 2 = 20$ \omterm{def:g7:atoms-and-molecules:molecule}{molekul} oksigen.

\textbf{11.} 10 \omterm{def:g7:atoms-and-molecules:molecule}{molekul} karbon dioksida.

\textbf{12.} $10 \times 2 = 20$ \omterm{def:g7:atoms-and-molecules:molecule}{molekul} air.
\end{solution}
